% !TEX root = ../main.tex

\section{Writing, Structure, and Navigation}
\label{sec:writing-navigation}

NeoSetzer continuously reads the source so that the editor can offer
completion, navigation, and structure information while you write. These
features are source-level helpers: they are designed to be fast and reliable,
not to emulate every possible TeX macro expansion.

\subsection{A \textit{nested} heading title}
\label{sec:writing-navigation:nested-title}

This heading intentionally contains a nested command. It is a useful sample
for the Document Structure sidebar: the title remains visible as one heading,
while its italic formatting belongs to the rendered PDF. Try selecting it from
the sidebar and compare the source, outline, and preview.

\subsection{Labels, references, and completion}
\label{sec:writing-navigation:labels}

This section has a label, so another part of the document can refer to
Section~\ref{sec:writing-navigation}. NeoSetzer indexes labels, references,
citations, and bibliography keys from open documents to improve completion.
Type a backslash for command completion, or use the configurable manual
completion shortcut when you prefer to invoke the list explicitly.

\subsection{TODO items and comments}
\label{sec:writing-navigation:todos}

The following visible note is also a small navigation sample:
\todo{Replace this note after exploring the TODO list in Document Structure.}

NeoSetzer recognizes \texttt{\textbackslash todo\{...\}} items and lists them
in the document-structure sidebar. Select the TODO entry to jump to its source
location. In a real project you can delete the command after finishing the
task; this example keeps it as a discoverable demonstration.

\subsection{Open a child file directly}
\label{sec:writing-navigation:root}

The first line of this file contains \texttt{\% !TEX root = ../main.tex}.
Open this child file in NeoSetzer and build: the relative root declaration
points the build at \texttt{main.tex}, while keeping the child convenient to
edit. This pattern is useful for reports, theses, and books split into chapters.
